\section{Which Cache Information Changes Allocation?}
\label{sec:graph-information}
The allowed query demands determine the cost information a search must retain.
A cache is sufficient for restoration only when the compatible execution objects
are available. For pricing, several queries can sometimes be aggregated.

\begin{theorem}[Demand-profile quotient]
\label{thm:demand-quotient}
Fix the family of allowed panels $\mathcal Q_1,\ldots,\mathcal Q_m$.
Group queries by their membership profiles
$q\sim r$ if $\ind\{q\in\mathcal Q_j\}=\ind\{r\in\mathcal Q_j\}$ for every $j$.
For each group $C$ and prefix $u$, store
\begin{equation}
 R_K(C,u)=\sum_{q\in C}\ind\{(q,u)\notin K\}w_{q,u}.
 \label{eq:group-residual}
\end{equation}
The bill for an allowed demand $\mathcal Q_j\times P_x$ equals
\[
 \sum_{u\in P_x}\sum_{C\subseteq\mathcal Q_j}R_K(C,u).
\]
Publishing that demand sets precisely these group--prefix cells to zero.
Conversely, merging two queries that some allowed panel separates loses the
ability to recover every panel bill from group residual sums, already with
unit weights and arbitrary binary residual states.
\end{theorem}

\paragraph{Proof.}
Every allowed panel contains either all or none of a profile group. Partitioning
the query sum therefore gives its exact bill, and publication removes precisely
the demanded cells. For the converse, swap one unit of uncached demand between
two merged queries. The group total stays fixed, while a separating panel's bill
changes. A common summary cannot recover both bills. Lean checks both directions.

\begin{proposition}[Full-panel compression]
\label{prop:graph-compression}
Fix the panel $\mathcal Q$ and resource weights. For an arbitrary cache $K$, define the residual prefix bill
\begin{equation}
 r_K(u)=\sum_{q\in\mathcal Q}\ind\{(q,u)\notin K\}w_{q,u}.
 \label{eq:residual-prefix}
\end{equation}
For a full-panel demand $\mathcal Q\times P_x$, its cold bill is $\sum_{u\in P_x}r_K(u)$. Publishing this demand sets $r_K(u)$ to zero on $P_x$ and leaves the other residual bills unchanged. If $K=\mathcal Q\times S$, the bill further reduces to $\sum_{u\in P_x\setminus S}\bar w_u$, where $\bar w_u=\sum_qw_{q,u}$. Empty initialization followed by full-panel publication preserves this rectangular form.
\end{proposition}

The main experiment uses this fixed-full-panel case. Its completion bills admit
an exact prefix-level quotient. The qid--prefix model also specifies the physical
checkpoints and the query-conditioned price observations; these are separate
from the aggregated cost totals. Online price revisions refresh the residual
sums, as detailed in the supplementary material.

Selective demands can require finer information. Consider two queries, prefixes
$u,v$, unit prices, and the caches
\[
 K_1=\{(q_0,u),(q_1,v)\},\quad
 K_2=\{(q_1,u),(q_0,v)\}.
\]
Both prefixes have one full-panel residual unit in either state.
On query $q_0$, their bills are $(0,1)$ and $(1,0)$. Equal-quality allocation
therefore reverses. The supplement gives a separation reachable from empty
caches through complete overlapping-panel evaluations. The result identifies
necessary information within partition-based residual summaries; bitsets and
other lossless encodings can preserve that same information.
